% !TeX root = ../main.tex

\chapter{Analisi Economica: SLO, Costi e ROA}
\label{cap:analisi_economica}

\section{SLO e Costo del Downtime}
\label{sec:slo_downtime}

Sono stati definiti tre SLO coerenti col workload:
\begin{itemize}
    \item \textbf{SLO 1 -- Latenza assegnazione:} 90\% degli ordini \textit{high priority} assegnati entro 1 minuto (giustificato dal \texttt{time.sleep(30)} + latenza inferenza).
    \item \textbf{SLO 2 -- Disponibilità data layer:} 99{,}9\% di uptime per MQTT e InfluxDB; una perdita prolungata renderebbe cieco l'\texttt{HealthAgent}.
    \item \textbf{SLO 3 -- Affidabilità agentica:} 99\% di tool call valide senza innesco di \texttt{MAX\_AGENT\_ITERATIONS}.
\end{itemize}

Nel dominio della logistica critica, un'ora di downtime è stimata a circa \textbf{15.000\,\texteuro/ora} (media dei guadagni orari + penali contrattuali per SLA violati).

\section{Return on Agent (ROA): OPEX e CAPEX}
\label{sec:roa}

L'orchestrazione autonoma introduce un bilanciamento critico tra i costi operativi delle API (OPEX) e il risparmio di capitale hardware (CAPEX) ottenuto via logiche predittive.

\textbf{Ottimizzazione OPEX.} Il sistema limita l'impatto economico delle chiamate LLM attraverso tre strategie integrate:
\begin{itemize}
    \item \textit{Triage preventivo}: il \texttt{triage\_manager} analizza lo snapshot del sistema prima di invocare i modelli, attivando \texttt{LogisticAgent} ed \texttt{HealthAgent} solo se sussistono reali condizioni di sbilanciamento (es. ordini pendenti accoppiati a droni idle, o necessità di scaling/manutenzione), azzerando i costi dei cicli no-op ridondanti.
    \item \textit{Ingegneria dei prompt}: i prompt di sistema impongono restrizioni sintattiche che vietano la generazione di testo libero, preamboli o spiegazioni discorsive; l'output è confinato alle sole chiamate di funzione (\textit{tool calling}), minimizzando i token di completamento.
    \item \textit{Monitoraggio consumi}: auditing sistematico dei token (prompt, completamento, totale) al termine di ogni esecuzione, per rendicontazione finanziaria precisa del processo decisionale.
\end{itemize}

\textbf{Preservazione CAPEX.} Il ritorno sull'investimento si concretizza nella mitigazione dei rischi d'impresa e nell'ottimizzazione infrastrutturale:
\begin{itemize}
    \item \textit{Manutenzione preventiva dinamica}: il \texttt{LogisticAgent} include nei prompt la regola di evitare carichi $>3$\,kg su droni con usura $>70\%$ (\textit{soft constraint} prompt-driven), riducendo il rischio di esaurimento batteria e perdita del vettore.
    \item \textit{Elasticità infrastrutturale cloud}: l'\texttt{HealthAgent} modula dinamicamente le repliche del Deployment in base al carico effettivo, riducendo i costi di over-provisioning.
    \item \textit{Operational Safety Cap}: limite massimo di 6 droni in autonomia (\texttt{POLICY\_LIMITS}); qualsiasi espansione ulteriore richiede \texttt{request\_human\_approval}, ponendo un \textit{budget cap} rigido contro loop di scaling incontrollati.
    \item \textit{Massimizzazione del throughput}: la coda ordini è ordinata per priorità, con l'agente che alloca la capacità residua privilegiando i task \textit{high priority}.
\end{itemize}

\subsection*{Proiezioni Economiche}
Ipotizzando 2.880 cicli/giorno (polling ogni 30\,s):

\begin{table}[ht]
\centering
\begin{tabular}{lrr}
\toprule
\textbf{Voce} & \textbf{Gemini 2.5 Pro} & \textbf{Mistral Small} \\
\midrule
Infrastruttura cloud (1 control-plane + 2 worker) & 150{,}00 \$/mese & 150{,}00 \$/mese \\
Token LLM & $\sim$500{,}00 \$/mese & $\sim$10{,}00 \$/mese \\
\midrule
\textbf{OPEX totale mensile} & \textbf{$\sim$650{,}00 \$/mese} & \textbf{$\sim$160{,}00 \$/mese} \\
\bottomrule
\end{tabular}
\caption{Confronto OPEX al variare del modello LLM.}
\label{tab:costi_opex}
\end{table}

La forbice di costo (delta $\sim$490\,\$/mese) consente di modulare la scelta del modello: \textit{Mistral Small} per scenari ordinari di instradamento; \textit{Gemini 2.5 Pro} riservato a finestre critiche con anomalie complesse o conflitti di policy.

\section{Trade-off Costo/Affidabilità e Automazione/Sicurezza}
\label{sec:analisi_tradeoff}

\textbf{Costo vs.\ affidabilità.} Un sistema deterministico presenta OPEX minimi (esecuzione di script locali) ma soffre di scarsa resilienza: a fronte di eccezioni non gestite o anomalie di rete subisce arresti critici, e il ripristino richiede l'intervento manuale, incrementando i costi di reperibilità oraria e introducendo latenze di risoluzione. Il sistema multi-agente paga un costo variabile in token (fino a 500\,\$/mese con \textit{Gemini 2.5 Pro}) in cambio della flessibilità semantica del modello, che contestualizza l'errore, ragiona sulla deviazione e devia su flussi alternativi senza richiedere intervento umano.

\textbf{Automazione vs.\ sicurezza.} Il deterministico garantisce l'osservanza rigorosa delle policy ma è rigido davanti ad anomalie inattese (es. sovrapposizione di backlog di ordini e fermo contemporaneo di più vettori), provocando blocchi fino all'intervento manuale. Il multi-agente massimizza l'efficienza decisionale abilitando valutazioni di trade-off complessi tra metriche di business e vincoli fisici; l'elasticità decisionale introduce però rischi (allucinazioni nelle tool call, derive logiche sullo scaling) mitigati dai guardrail del Capitolo~\ref{cap:safety_guardrails}.
